\section{Preliminaries}\label{sec:preliminaries}


%\subsection{Preliminaries}
%We establish the concept of {time} within a domain denoted as \emph{Time}, where the instances of time possess a total order denoted by $\leq$. The general definitions presented hereforth remain unaffected whether time is considered discrete (\(\mathbb{N}^0\)) or continuous (\(\mathbb{R}_{\geq 0}\)). A segment of the time domain refers to a subset $T$ that is closed downward, meaning that if $t \in T$, then all elements $t'$ such that $t'\leq t$ also belong to $T$. 


The framework introduced in Section~\ref{sec:architecture} relies on formal models to capture PT and DT views, reason about their behavior, and establish alignment. We now introduce the formal foundations underlying these models: timed transition systems as the semantic domain, timed automata as their syntactic representation, and weak timed bisimulation as the behavioral equivalence relation on which alignment will be grounded.

We model PT and DT views as \emph{Timed Transition Systems (TTS)}~\cite{ALUR1994183, Baier08}, the standard semantic model for real-time systems capturing both discrete event-driven behavior and the passage of time. A TTS is defined as follows.

%\ {Timed Transition Systems (TTS)} is the standard semantic model for real-time systems 
%\cite{ALUR1994183, Baier08} 
%that captures both discrete event-driven behavior and the passage of time, allowing us to reason about the temporal aspects of system behavior. Formally, a TTS is defined as follows:
%
%In order to model a system and its components, we utilize the concept of timed transition systems (TTS). A TTS captures both the discrete transitions between states and t he passage of time, allowing us to reason about the temporal aspects of system behavior. Formally, a TTS is defined as follows:




\begin{definition}%[Timed Transition System]
  \label{def:tts}
A {Timed Transition System} is a tuple
$\mathcal{T} = (S, \discreteAlphabet, \delayAlphabet, \rightarrow)$, where:
\begin{itemize}
    \item $S$ is a set of states%, where each state $s \in S$ is a tuple of the form \( s = (l, v) \), with  \( l \) a discrete location and \( v \) a valuation of clocks, mapping each clock to a real value.
    \item $\discreteAlphabet$ is a set of \emph{event labels},
    \item $\delayAlphabet$ is a set of \emph{delay labels}, with $\delta$ representing the elapse of $\delta$ time units. 
    
    
    %δ∈ΣD​ represents the elapse of δ\delta δ time units."
    \item $\rightarrow \subseteq S \times (\discreteAlphabet \cup \delayAlphabet) \times S$
    is the transition relation.
\end{itemize}
\end{definition}
%A sequence of events and delays induced by an execution of a TTS is called a trace. The semantics of $\comp$, denoted by $\sem{\comp}$, is defined as the set of all traces induced by its TTS.









We make use of \emph{Timed Automaton (TA)} \cite{ALUR19932, Bouyer2005AnIT} as syntactic representation of each model. A TA is defined as follows. %A TA extends finite automata with real-valued clocks to capture quantitative timing constraints. Formally, 

\begin{definition}%[Timed Automaton]
  \label{def:ta}
    A Timed Automaton is a tuple
\(
TA = (L, \ell_0, C, \Sigma, E, \text{Inv}),
\) 
where:
\begin{itemize}
  \item $L$ is a finite set of locations with initial location $\ell_0 \in L$,
  \item $C$ is a finite set of clocks,
  \item $\Sigma = \discreteAlphabet \cup \delayAlphabet$ is the action alphabet, partitioned into event-driven and delay-driven actions with $\discreteAlphabet \cap \delayAlphabet = \emptyset$,
  \item $E \subseteq L \times \mathcal{B}(C) \times \Sigma \times 2^C \times L$ is a finite set of edges, where $\mathcal{B}(C)$ denotes clock constraints,
  \item $\text{Inv}: L \to \mathcal{B}(C)$ assigns invariants restricting the admissible time elapse in each location.
\end{itemize}
\end{definition}






The semantics of a TA is given as a TTS as of Definition~\ref{def:tts}, defined as of \cite{ALUR19932, Bouyer2005AnIT}, and written as $\sem{TA}$. A \emph{state} of the corresponding TTS is a pair $(\ell, v) \in L \times \mathbb{R}_{\geq 0}^C$, where $\ell$ is the current location and $v$ a clock valuation. %A state is valid only if $v \models \text{Inv}(\ell)$. 
Discrete transitions correspond to edges in $E$ with clock resets, while delay transitions correspond to uniform advancement of all clocks under the invariant of the current location.  

To formalize the expected behavior of systems, we use Timed Computation Tree Logic (TCTL) properties \cite{ALUR1994183}. The semantics of a property $\phi$, denoted by $\sem{\phi}$, is the set of traces that satisfy $\phi$. A TA $TA$ satisfies a formula $\phi$, written $TA \models \phi$, iff $\sem{TA} \subseteq \sem{\phi}$. 
A property $\varphi$ over an alphabet $X$ is built from labels in $X$ using boolean connectives ($\land$, $\lor$, $\neg$) and timed path quantifiers ($\mathbf{A}$, $\mathbf{E}$, $\mathbf{G}$, $\mathbf{F}$, $\mathbf{U}$); we refer to~\cite{ALUR1994183} for the full syntax. Given a function $\sigma : X \to Y$, we write $\sigma(\varphi)$ for the property obtained by replacing each label $a$ in $\varphi$ with $\sigma(a)$.
%A contract $C = (A, G)$ is a pair, where $A$ and $G$ are sets of traces representing, respectively,
%the assumptions on the environment and the guarantees provided by the component. The semantics of a contract $C$, denoted by $\sem{C}$, is defined as \( \sem{C} \triangleq \{\, \sigma \mid \sigma \in A \implies \sigma \in G \,\} \).A component $\comp$ {satisfies} a contract $C$, written $\comp \models C$, if $\sem{\comp} \subseteq \sem{C}$
%

%the set of traces for which the assumptions imply the guarantees, formally \( \sem{C} \triangleq \{\, \sigma \mid \sigma \in A \implies \sigma \in G \,\} \). A component $\comp$ {satisfies} a contract $C$, written $\comp \models C$, if $\sem{\comp} \subseteq \sem{C}$. %A contract $C' = (A', G')$ refines another contract $C = (A, G)$, written $C' \refines C$, if $A \subseteq A'$ and $G' \subseteq G$. This ensures that any component satisfying $C'$ also satisfies $C$. 

%However, in systems with hard real-time and cyclic behaviors, these infinite paths can be represented finitely by identifying repeating patterns  \cite{Dou02}. 
%Verifying trace inclusion in timed systems is generally undecidable due to infinite execution paths leading to unbounded state spaces \cite{ALUR1994183}. To overcome this problem, we employ \emph{symbolic abstractions} based on zones to obtain a finite representation of the timed system \cite{ALUR19932, Bouyer2005AnIT}. A \emph{zone} is a convex set of clock valuations that satisfies the same set of constraints. A TA $A$ can thus be abstracted into a \emph{zone-based symbolic automaton} $\Gamma(A)$, whose states are pairs $(\ell, Z)$ with $\ell \in L$ and $Z$ a zone representing all valuations reachable at location $\ell$. Every concrete state $(\ell, v)$ belongs to exactly one symbolic state $(\ell, Z)$ with $v \in Z$. Transitions between zones are induced both by discrete edges and by the passage of time within each location. This symbolic construction yields a \emph{finite-state abstraction} of the TA that preserves all relevant properties. 


%A network $\mathcal{N} = \mathit{TA}_1 \parallel \mathit{TA}_2 \parallel \cdots \parallel \mathit{TA}_n$ is a parallel composition as of \cite{Uppaal} in which each $\mathit{TA}_i$ has local clocks $C_i$ (clock sets pairwise disjoint) and an action alphabet $\Sigma^{(i)}$. Actions in $\Sigma^{(i)} \setminus \bigcup_{j \ne i} \Sigma^{(j)}$ interleave freely, while actions shared between two or more components are {synchronizing} and must be taken simultaneously by all sharing components. Time advances globally and simultaneously across all components, subject to all location invariants. %We write $\mathit{Sync}(\mathit{TA}_i) = \Sigma^{(i)} \cap \bigcup_{j \ne i} \Sigma^{(j)}$ for the synchronization alphabet of $\mathit{TA}_i$ within $\mathcal{N}$. 
%The TTS semantics of a network is defined analogously to that of a single TA, with discrete transitions respecting the handshake discipline and time advancing jointly.



To reason about the behavioral relationship between two TTS, we 
use weak timed simulation (WTS) and weak timed bisimulation (WTB). Intuitively, 
a simulation captures the idea that one system can match every 
observable move of another, while abstracting over internal 
transitions. Bisimulation strengthens 
this to a symmetric requirement: neither system can exhibit an 
observable behaviour the other cannot match.

Formally, internal actions are modeled as $\tau$-transitions, 
which arise from projection of event labels outside the alphabet 
of interest. We write $s \xRightarrow{\varepsilon} s'$ if $s'$ is 
reachable from $s$ by any finite sequence of $\tau$-transitions 
and delay transitions, and $s \xRightarrow{a} s'$ for $a \in 
\Sigma_E \setminus \{\tau\}$ if there exist $s_1, s_2$ such that 
$s \xRightarrow{\varepsilon} s_1 \xrightarrow{a} s_2 
\xRightarrow{\varepsilon} s'$.

\begin{definition}
A relation $R \subseteq S_1 \times S_2$ is a WTS 
from $\mathcal{T}_1$ to $\mathcal{T}_2$ on $X \subseteq 
\Sigma_{E_1} \cap \Sigma_{E_2}$ iff $(s^1_0, s^2_0) \in R$ and for every $(s, t) \in R$:
\begin{itemize}
\item $\forall a \in X,s'$ s.t. $s \xRightarrow{a} s'$, 
     $\exists t'$ with $t \xRightarrow{a} t'$ and 
    $(s', t') \in R$;
\item $\forall \delta \in \Sigma_D$ and $s'$ with $s 
\xrightarrow{\delta} s'$,  $\exists t'$ reachable from 
$t$ by $\xRightarrow{\varepsilon}$ with elapsed time $\delta$ 
    such that $(s', t') \in R$.
\end{itemize}
We write $\mathcal{T}_1 \refines{} \mathcal{T}_2$ iff such a 
relation exists.
\end{definition}

\begin{definition}
A relation $R$ is a WTB on $X$ iff 
both $R$ and $R^{-1}$ are WTSs on $X$. We write 
$\mathcal{T}_1 \aligns{} \mathcal{T}_2$.
\end{definition}

As a standard result, if $\mathcal{T}_1 \aligns{} \mathcal{T}_2$ on alphabet $X$, then for every TCTL property $\varphi$ over $X$: $\mathcal{T}_1 \models \varphi \iff \mathcal{T}_2 \models \varphi$ \cite{Baier08}. If $\mathcal{T}_1 \refines{} \mathcal{T}_2$ on alphabet $X$, then for every safety property $\varphi$ over $X$: $\mathcal{T}_2 \models \varphi \Rightarrow \mathcal{T}_1 \models \varphi$ \cite{Baier08}.

%
%Given a TA $A$ and a set of event labels $X \subseteq \discreteAlphabet$, the \emph{projection of $A$ to $X$}, written $\projectionX A$, is the TA obtained from $A$ by relabeling every edge with action $a \in \discreteAlphabet \setminus X$ as $\tau$, a distinguished silent event. Locations, clocks, guards, resets, invariants, and delay labels are unchanged. Projection lifts to the TTS semantics as $\sem{\projectionX A} = \projectionX{\sem{A}}$, where on TTSs $\projectionSkeletonX$ relabels every non-$X$ event transition as $\tau$. Projection is monotone (if $X \subseteq Y$ then $\projectionX A$ hides at least as much as $\projection{Y} A$) and idempotent ($\projectionX A \circ \projection{Y} A = \projection{X \cap Y} A$). After any projection, in any resulting TTS, we write $s \xRightarrow{\varepsilon} s'$ iff $s'$ is reached from $s$ by any (possibly empty) sequence of $\tau$-transitions and delay transitions. For $a \in \discreteAlphabet \setminus \{\tau\}$, write $s \xRightarrow{a} s'$ iff there exist $s_1, s_2$ with $s \xRightarrow{\varepsilon} s_1 \xrightarrow{a} s_2 \xRightarrow{\varepsilon} s'$.
%
%
%
%Let $A_1, A_2$ be TAs and $X \subseteq \discreteAlphabet^{(1)} \cap \discreteAlphabet^{(2)}$. We define the following relations:
%
%\begin{definition}%[$X$-weak simulation; refinement relation $\refines{X}$]
%\label{def:refinement}
%A relation $R \subseteq S_1 \times S_2$ is an \emph{$X$-weak simulation from $A_1$ to $A_2$} iff $(s_0^1, s_0^2) \in R$ and for every $(s,t) \in R$:
%\begin{itemize}
%  \item $\forall a \in X$ and $s'$ with $s \xrightarrow{a} s'$ in $\projection{X}{\sem{A_1}}$,$\exists t'$ with $t \xRightarrow{a} t'$ in $\projection{X}{\sem{A_2}}$ and $(s',t') \in R$;
%  \item  $\forall\delta \in \Sigma_D$ and $s'$ with $s \xrightarrow{\delta} s'$,  $\exists t'$ reached from $t$ by $\xRightarrow{\varepsilon}$ with elapsed time $\delta$, s. t. $(s',t') \in R$.
%\end{itemize}
%We write $A_1 \refines{X} A_2$ iff such a simulation exists.
%\end{definition}
%
%\begin{definition}%[$X$-weak bisimulation; alignment relation $\aligns{X}$]
%\label{def:alignment}
%A relation $R$ is an \emph{$X$-weak bisimulation} iff both $R$ and $R^{-1}$ are $X$-weak simulations. Write $A_1 \aligns{X} A_2$.
%\end{definition}
%

%The following properties hold by standard arguments and are used throughout:
%\begin{itemize}
  %\item {(Reflexivity)} $A \refines{} A$ and $A \aligns{} A$.
  %\item {(Transitivity)} $\refines{}$ and $\aligns{}$ are transitive.
  %\item {(Bisimulation implies simulation)} $A_1 \aligns{} A_2 \Rightarrow A_1 \refines{} A_2$ and $A_2 \refines{} A_1$.
  %\item {(Projection monotonicity)} $X \subseteq Y \Rightarrow (A_1 \refines{Y} A_2 \Rightarrow A_1 \refines{} A_2)$; analogously for $\aligns{}$ and $\aligns{}$.
  %\item {(TCTL preservation)} If $A_1 \aligns{} A_2$ on alphabet $X$, then for every TCTL propery $\varphi$ over $X$: $A_1 \models \varphi \iff A_2 \models \varphi$
  %\item {(Safety preservation)} If $A_1 \refines{} A_2$ on alphabet $X$, then for every safety property $\varphi$ over $X$: $A_2 \models \varphi \Rightarrow A_1 \models \varphi$.
%
  %\item {(Precongruence for $\parallel$)} If $A_1 \refines{} A_1'$, $A_2 \refines{} A_2'$, and the synchronization alphabet of the composed network lies in $X$, then $A_1 \parallel A_2 \refines{} A_1' \parallel A_2'$; analogously for $\aligns{}$.
%\end{itemize}
%
%
%\begin{itemize}
%  \item Remove the contracts etc, maybe, and add refinement in the form of simulation and bisimulation over projected alphabets. 
%\end{itemize}
%
%
%%Components are composed into systems via \emph{networks of timed automata} in the UPPAAL style \cite{Uppaal}. 
%

%We establish the concept of {time} within a domain denoted as \emph{Time}, where the instances of time possess a total order denoted by $\leq$. The general definitions presented hereforth remain unaffected whether time is considered discrete (\(\mathbb{N}^0\)) or continuous (\(\mathbb{R}_{\geq 0}\)). A segment of the time domain refers to a subset $T$ that is closed downward, meaning that if $t \in T$, then all elements $t'$ such that $t'\leq t$ also belong to $T$. 


%In order to model a system or its components, we utilize the concept of timed transition systems (TTS), which are a widely adopted formalism for representing the behavior of real-time systems. A TTS captures both the discrete transitions between states and the passage of time, allowing us to reason about the temporal aspects of system behavior. Formally, a TTS is defined as follows:
%
%
%
%\begin{definition}[Timed Transition System]\label{def:tts}
%A \emph{Timed Transition System (TTS)} is a tuple
%$\mathcal{T} = (S, \Sigma_E, \Sigma_D, \rightarrow)$, where:
%\begin{itemize}
%    \item $S$ is a set of states%, where each state $s \in S$ is a tuple of the form \( s = (l, v) \), with  \( l \) a discrete location and \( v \) a valuation of clocks, mapping each clock to a real value.
%    \item $\Sigma_E$ is a set of \emph{event labels},
%    \item $\Sigma_D$ is a set of \emph{delay labels}, with $\delta \in \Sigma_D$ representing the elapse of $\delta \in Time$ time units. 
%    
%    
%    %δ∈ΣD​ represents the elapse of δ\delta δ time units."
%    \item $\rightarrow \subseteq S \times (\Sigma_E \cup \Sigma_D) \times S$
%    is the transition relation.
%\end{itemize}
%\end{definition}
%A sequence of events and delays induced by an execution of a TTS is called a trace. The semantics of $\comp$, denoted by $\sem{\comp}$, is defined as the set of all traces induced by its TTS.
%
%
%
%To formalize the expected behavior of components, we use temporal logic formulas, such as Timed Computation Tree Logic (TCTL) \cite{ALUR1994183}, to define sets of acceptable traces. The semantics of a formula $\phi$, denoted by $\sem{\phi}$, is the set of traces that satisfy $\phi$. A component $\comp$ satisfies a formula $\phi$, written $\comp \models \phi$, iff $\sem{\comp} \subseteq \sem{\phi}$. 
%
%
%
%
%
%We make use of \emph{Timed Automaton (TA)} and \emph{Region Automaton} \cite{ALUR19932, Bouyer2005AnIT} as a syntactic representation of the component models. A TA extends finite automata with real-valued clocks to capture quantitative timing constraints. Formally, 
%
%\begin{definition}[Timed Automaton]\label{def:ta}
%    A Timed Automaton is a tuple
%\(
%TA = (L, \ell_0, C, \Sigma, E, \text{Inv}),
%\) 
%where:
%\begin{itemize}
%  \item $L$ is a finite set of locations with initial location $\ell_0 \in L$,
%  \item $C$ is a finite set of clocks,
%  \item $\Sigma = \Sigma_E \cup \Sigma_D$ is the action alphabet, partitioned into event-driven and delay-driven actions with $\Sigma_E \cap \Sigma_D = \emptyset$,
%  \item $E \subseteq L \times \mathcal{B}(C) \times \Sigma \times 2^C \times L$ is a finite set of edges, where $\mathcal{B}(C)$ denotes clock constraints,
%  \item $\text{Inv}: L \to \mathcal{B}(C)$ assigns invariants restricting the admissible time elapse in each location.
%\end{itemize}
%\end{definition}
%
%
%
%
%
%
%The semantics of a TA is given as a TTS as of Definition~\ref{def:tts}. A \emph{state} of the corresponding TTS is a pair $(\ell, v) \in L \times T^C$, where $\ell$ is the current location and $v$ a clock valuation. %A state is valid only if $v \models \text{Inv}(\ell)$. 
%Discrete transitions correspond to edges in $E$ with clock resets, while delay transitions correspond to uniform advancement of all clocks under the invariant of the current location.  
%
%%However, in systems with hard real-time and cyclic behaviors, these infinite paths can be represented finitely by identifying repeating patterns  \cite{Dou02}. 
%Verifying trace inclusion in timed systems is generally undecidable due to infinite execution paths leading to unbounded state spaces \cite{ALUR1994183}. To overcome this problem, we employ \emph{symbolic abstractions} based on zones to obtain a finite representation of the timed system \cite{ALUR19932, Bouyer2005AnIT}. A \emph{zone} is a convex set of clock valuations that satisfies the same set of constraints. A TA $A$ can thus be abstracted into a \emph{zone-based symbolic automaton} $\Gamma(A)$, whose states are pairs $(\ell, Z)$ with $\ell \in L$ and $Z$ a zone representing all valuations reachable at location $\ell$. Every concrete state $(\ell, v)$ belongs to exactly one symbolic state $(\ell, Z)$ with $v \in Z$. Transitions between zones are induced both by discrete edges and by the passage of time within each location. This symbolic construction, first introduced in \cite{ALUR19932} and extended in \cite{Bouyer2005AnIT}, yields a \emph{finite-state abstraction} of the TA that preserves all relevant properties. 